\documentclass[a4paper,12pt]{article}
\usepackage[utf8]{inputenc}
\usepackage{geometry}
\usepackage{graphicx}
\usepackage{listings}
\usepackage{xcolor}
\usepackage{float}
\usepackage{hyperref}

% Config
\definecolor{codegreen}{rgb}{0,0.6,0}
\definecolor{codegray}{rgb}{0.5,0.5,0.5}
\definecolor{backcolour}{rgb}{0.95,0.95,0.92}

\lstdefinestyle{mystyle}{
    backgroundcolor=\color{backcolour},   
    commentstyle=\color{codegreen},
    keywordstyle=\color{magenta},
    basicstyle=\ttfamily\footnotesize,
    breakatwhitespace=false,         
    breaklines=true,                 
    captionpos=b,                    
    keepspaces=true,                 
    numbers=left,                    
    numbersep=5pt,                  
    showspaces=false,                
    showtabs=false,                  
    tabsize=2,
    frame=single
}
\lstset{style=mystyle}

\title{Practical Work 6: GlusterFS}
\author{Le Quy Ton \\ 23BI14425}
\date{\today}

\begin{document}

\maketitle

\section{Introduction}
GlusterFS is a scalable network-attached storage file system. This practical work aims to install GlusterFS, configure a distributed replicated volume, and evaluate its performance through benchmarking using Python automation.

\section{Setup and Configuration}

\subsection{Installation}
We installed GlusterFS server on our Debian environment using the following commands:
\begin{lstlisting}[language=bash]
sudo apt-get update
sudo apt-get install -y glusterfs-server
sudo systemctl start glusterd
sudo systemctl enable glusterd
\end{lstlisting}

\subsection{Trusted Pool and Volume Creation}
We used a Python automation script to detect the machine's real IP address (as GlusterFS prohibits `localhost` for volume creation) and configured a \textbf{Distributed Replicated Volume} (Replica count = 2).

\textbf{1. Prepare Bricks:}
\begin{lstlisting}[language=bash]
mkdir -p /data/brick1/gv0
mkdir -p /data/brick2/gv0
\end{lstlisting}

\textbf{2. Create Volume:}
\begin{lstlisting}[language=bash]
# Using the machine's Real IP (detected automatically)
sudo gluster volume create gv0 replica 2 \
    <Host_IP>:/data/brick1/gv0 \
    <Host_IP>:/data/brick2/gv0 force
\end{lstlisting}

\textbf{3. Start and Mount Volume:}
\begin{lstlisting}[language=bash]
sudo gluster volume start gv0
mkdir -p /mnt/gluster_client
# Mounting using the Real IP
sudo mount -t glusterfs <Host_IP>:/gv0 /mnt/gluster_client
\end{lstlisting}

\section{Benchmarking}
We utilized a Python script to automate the benchmarking process. The tests focused on two scenarios required by the specifications.

\subsection{Methodology}
\begin{itemize}
    \item \textbf{Large Files:} We generated a 512MB file filled with zeros. We cleared the RAM cache (\texttt{drop\_caches}) before reading to ensure we measured disk I/O, not RAM speed. The metric is Read Speed (MB/s).
    \item \textbf{Small Files:} We created 500 small empty files sequentially in the GlusterFS mount point. The metric is Number of Accesses (Creations) per second.
\end{itemize}

\subsection{Results}
The benchmarks were performed on a single-node setup (simulating 2 bricks). The results obtained from the automation script are shown below.

\begin{table}[H]
    \centering
    \begin{tabular}{|c|c|c|}
        \hline
        \textbf{Configuration} & \textbf{Small Files (ops/s)} & \textbf{Large Files Read (MB/s)} \\
        \hline
        1 Server (2 Bricks) & 190.28 & 215.17 \\ 
        \hline
        \textit{Hypothetical 2 Servers} & \textit{Estimated ~160} & \textit{Estimated ~400} \\
        \hline
    \end{tabular}
    \caption{Performance Benchmarks}
    \label{tab:benchmark}
\end{table}

\begin{figure}[H]
    \centering
    \includegraphics[width=1.0\linewidth]{image1.png}
    \caption{Successful execution and benchmark results}
    \label{fig:execution}
\end{figure}

\section{Conclusion}
We successfully set up a GlusterFS environment with a distributed replicated volume using an automated Python script to handle IP detection and configuration.
\begin{itemize}
    \item \textbf{Replication} ensures data availability; if one brick fails, data is still accessible.
    \item \textbf{Performance}: We achieved a read speed of \textbf{215.17 MB/s} for large files and a creation rate of \textbf{190.28 ops/s} for small files. The overhead for small files is expected due to metadata synchronization in a distributed file system.
\end{itemize}

\end{document}
